\section*{Asesmen}
\addcontentsline{toc}{section}{Asesmen}

Asesmen Bab IX mencakup Praktikum 2 (12\%) dan UTS Praktik (14\%).
Fokusnya adalah kemampuan menguji grammar, men-\textit{debug} konflik,
dan menghasilkan berkas grammar yang bebas konflik.

\begin{table}[H]
\centering
\caption{Instrumen asesmen Bab IX}
\label{tab:ch09-asesmen}
\begin{tabularx}{\textwidth}{@{}lXl@{}}
\toprule
\textbf{Instrumen} & \textbf{Indikator} & \textbf{Bobot} \\
\midrule
Praktikum 2 & Berkas grammar bebas konflik, log debug sebelum/sesudah,
dan uji lima string & 12\% (bagian) \\
UTS Praktik & Uji dan debug ambiguitas, left recursion, dan left
factoring & 14\% (bagian) \\
Kuis Konsep & Ambiguitas, left recursion, dan left factoring &
formatif \\
\bottomrule
\end{tabularx}
\end{table}

Kriteria kelulusan: mahasiswa dapat mengidentifikasi konflik,
men-\textit{debug} grammar sampai bebas konflik, membuktikan kesetaraan
bahasa dengan lima string uji, dan menjelaskan aturan resolusi yang
dipakai.
